
\documentclass[12pt]{article}

\usepackage{amsmath}
\usepackage{amssymb}
\usepackage{graphicx}
\usepackage{geometry}

\geometry{margin=1in}

\setlength{\parindent}{0pt}
\setlength{\parskip}{6pt}

\numberwithin{equation}{section}

\title{Physics Documentation}
\author{Thomas Walsh}
\date{\today}

\begin{document}

\maketitle
\tableofcontents
\newpage

\section{Introduction}

The objective of this project is to develop a computational 
model for rendering black holes.

In particular, this document is intended to translate my 
notes on relevant physical equations into a readable 
document on GitHub.

Additional documentation can be found within the 
\texttt{/Docs} directory of the repository.

\newpage

\section{Mathematical Foundations}

\subsection{Schwarzschild Radius and Null Geodesics}

The Schwarzschild radius is defined as

\begin{equation}
    r_s = \frac{2GM}{c^2}
    \label{eq:schwarzschild_radius}
\end{equation}

where $G$ is the gravitational constant, 
$M$ is the mass of the black hole, and 
$c$ is the speed of light.

We introduce the reciprocal radial coordinate $u$, defined as

\begin{equation}
    u = \frac{1}{r}
\end{equation}

Using this substitution, the Schwarzschild null geodesic 
orbit equation can be expressed in geometrized units 
($G = c = 1$) as

\begin{equation}
    \frac{d^2u}{d\phi^2} + u = 3Mu^2
    \label{eq:orbit}
\end{equation}

where $\phi$ is the angular coordinate in the plane of 
the photon's trajectory. This equation assumes a 
non-rotating, uncharged black hole and a photon with 
nonzero angular momentum.

\subsection{Flat Spacetime}

With a mass of zero ($M = 0$), we get the simplified 
equation

\begin{equation}
    u'' + u = 0
\end{equation}

\begin{equation}
    u'' = -u
\end{equation}

In order to satisfy this equation, we use both sine and 
cosine as independent solutions. Since the differential 
equation is both linear and homogeneous, we take a linear 
combination of our independent solutions.

\begin{equation}
    u(\phi) = A\cos(\phi) + B\sin(\phi)
\end{equation}

where $A$ and $B$ are arbitrary constants determined by 
the initial state of the photon.

Using the reciprocal radial coordinate and the 
polar-to-Cartesian coordinate transformations, we obtain

\begin{equation}
    1 = Ar\cos(\phi) + Br\sin(\phi)
\end{equation}

\begin{equation}
    Ax + By = 1
\end{equation}

This demonstrates that the photon follows a straight-line 
trajectory in flat spacetime.

\subsection{Characteristic Schwarzschild Quantities}

In geometrized units ($G = c = 1$), equation 
\eqref{eq:schwarzschild_radius} becomes

\begin{equation}
    r_s = 2M
\end{equation}

For this implementation, we normalize the mass of the 
black hole to one ($M = 1$), giving us

\begin{equation}
    r_s = 2
\end{equation}

Next, we utilize known equations in Schwarzschild 
spacetime to define

\begin{gather}
    \text{Event Horizon:} \quad r_s = 2M = 2 \\
    \text{Photon Sphere:} \quad r_p = 3M = 3 \\
    \text{Critical Impact Parameter:} \quad
    b_c = 3\sqrt{3}M = 3\sqrt{3}
\end{gather}

The event horizon defines the boundary from within which 
light cannot escape to infinity.

The photon sphere defines the radius at which photons 
can maintain unstable circular orbits.

The critical impact parameter defines the threshold 
between photon capture and escape for photons incoming 
from infinity.

\newpage

\section{Numerical Methods}

\subsection{Conversion to a First-Order ODE System}

We begin by introducing new variables to reduce equation 
\eqref{eq:orbit} into a system of first-order ordinary 
differential equations.

\begin{equation}
    y_1 = u
    \qquad
    y_2 = \frac{du}{d\phi}
\end{equation}

From these definitions, we obtain

\begin{equation}
    \frac{dy_1}{d\phi} = y_2
\end{equation}

\begin{equation}
    \frac{dy_2}{d\phi} = 3My_1^2 - y_1
\end{equation}

We can express this system using a state vector

\begin{equation}
    \mathbf{y} =
    \begin{pmatrix}
        y_1 \\
        y_2
    \end{pmatrix}
\end{equation}

giving us

\begin{equation}
    \frac{d\mathbf{y}}{d\phi}
    =
    \mathbf{f}(\mathbf{y})
    =
    \begin{pmatrix}
        y_2 \\
        3My_1^2 - y_1
    \end{pmatrix}
\end{equation}

This is the first-order ODE system that will be 
numerically integrated.

\subsection{Fourth-Order Runge-Kutta (RK4)}

The fourth-order Runge-Kutta method approximates the 
solution of a differential equation by evaluating its 
derivative at four stages within a given step. These 
evaluations are then combined using a weighted average 
to approximate the next state of the system.

Let $h = \Delta\phi$ denote the angular step size.

\begin{equation}
    \mathbf{y}_{n+1} =
    \mathbf{y}_n +
    \frac{h}{6}
    (\mathbf{k}_1 + 2\mathbf{k}_2 +
    2\mathbf{k}_3 + \mathbf{k}_4)
\end{equation}

The $\mathbf{k}$ terms are defined as

\begin{gather}
    \mathbf{k}_1 =
    \mathbf{f}(\mathbf{y}_n) \\
    \mathbf{k}_2 =
    \mathbf{f}\left(
        \mathbf{y}_n + \frac{h}{2}\mathbf{k}_1
    \right) \\
    \mathbf{k}_3 =
    \mathbf{f}\left(
        \mathbf{y}_n + \frac{h}{2}\mathbf{k}_2
    \right) \\
    \mathbf{k}_4 =
    \mathbf{f}\left(
        \mathbf{y}_n + h\mathbf{k}_3
    \right)
\end{gather}

We also advance our angular coordinate

\begin{equation}
    \phi_{n+1} = \phi_n + h
\end{equation}

\subsection{Justification}

It is important to note that there are several numerical 
methods that we could have chosen to work with in 
Section 3.2. However, methods such as Euler's method 
and second-order Runge-Kutta (RK2) generally provide 
lower accuracy than RK4 for sufficiently small step 
sizes.

Under suitable smoothness conditions, the global 
truncation errors of these methods are

\begin{gather}
    \text{Euler's Method:} \quad O(h) \\
    \text{RK2:} \quad O(h^2) \\
    \text{RK4:} \quad O(h^4)
\end{gather}

RK4 provides a balance between computational cost and 
numerical accuracy. However, the angular step size must 
still be chosen carefully to maintain numerical accuracy 
and stability, particularly for trajectories near the 
critical impact parameter.

\newpage

\section{Numerical Results}

\newpage

\section{Integration with Ray Tracer}

\end{document}
